\section{Beyond the reporting boundary}\label{sec:boundary}

A representation of current state needs new observations. A reusable rule
usually needs less frequent revision, but still has a context and can be
invalidated. This is the useful distinction between the plotting table and the
portable pattern in Figure~\ref{fig:room-and-machine}; neither is exempt from
checking its basis.

Reporting also has a participation cost. An organisation may require updates
from its members; a voluntary collaboration must give participants a reason
and an affordable way to supply them. This motivates a design question, not a
measured market boundary: can publication and confirmation sustain a useful
record where compulsory reporting is unavailable?

\subsection*{Attribution without inventing agreement}

The working-tree provenance design joins tool-edit witnesses, file observations,
commit links and mission-clock lineage. Each edge retains its basis.
A clock records the declared working context; a path or retrieval match can
suggest another association. An explicit confirmation can strengthen that
attribution. The inferred association must remain distinguishable from the
confirming act.

``No mission'' is a valid result. A later claim should identify claimant and
basis, retain any competing claims, and preserve the fact that the work was
unclaimed at commit time. Such records can reveal work outside planned
missions without rewriting its history to make the plan appear complete.
Accept, decline and no-response must remain distinct here for the same reason
as in the operator-feedback model.

The dated 2026-08-24 probe demonstrated only part of this design: 343 file
observations and nine witnesses across fourteen roots, with eighty paths
unattributed and no delivered followups. The attribution namespace was not
loaded in that process. These are historical limits of that probe, not a
current verdict on the service. A new claim of operation needs a new witnessed
join and delivery.

Absence requires its own coverage evidence. A join with no result cannot show
that no work occurred unless its producers, population and observation window
are known. A heartbeat cannot establish task progress, and frequent retrieval
cannot establish that every eligible turn was observed. This is where the
apparatus pattern \emph{Monitors measure the work} applies.

\subsection*{Policy scores and realised outcomes}

Expected free energy scores predicted outcomes against declared preferences,
with epistemic terms supplied by the model. A realised payoff answers a
different question. Calling both a policy score does not establish that they
are the same quantity.

The working notes' four tests, \textsc{s-g1}--\textsc{s-g4}, ask whether a
score derives from predicted outcomes and preferences; whether its trajectory
can contain different actions; whether changing pattern wiring can change the
score; and whether stipulated components are disclosed. These tests identify
what a proposed policy demonstration must exhibit. A repeated-action rollout
can be a restricted policy, but cannot alone demonstrate sensitivity to the
ordering of different actions.

Appendix~\ref{app:snatch} demonstrates one of these distinctions concretely:
changing the consultation order of unchanged patterns changes a game's realised
payoff from $+3$ to $-5$. That is evidence that wiring matters in the experiment.
The payoff comparison alone does not establish all four tests for the War
Machine's expected-free-energy calculation, nor validate the live system.
The model's observations, preference carrier and actual consumers must be
checked separately; Section~\ref{sec:c-vector} describes a specific gap in that
composition.

Part~IV therefore proposes an attribution and participation discipline, with
bounded component evidence. It establishes neither general organisational
transfer nor an integrated policy-learning result.
